\section{Architecture et diagrammes UML}
\label{sec:archi}

L'ensemble des diagrammes UML est rédigé en syntaxe \textbf{Mermaid} dans
le fichier \texttt{docs/diagrams.md} et exporté en PNG via la CLI
\texttt{mmdc} pour intégration dans le présent rapport. Cette section
résume les diagrammes les plus importants~; les versions complètes (et
zoomables) sont disponibles dans le dépôt.

\subsection{Vue d'ensemble (diagramme de composants)}

L'application suit une architecture web classique en 3-tiers. Le navigateur
échange exclusivement avec le serveur Django, qui orchestre les accès à la
base PostgreSQL, au système de fichiers (créatifs) et à l'API externe
DeepSeek.

\begin{itemize}
    \item \textbf{Couche présentation~:} templates Django + Tailwind +
    HTMX/Alpine. Pas de SPA, pas de bundler JS.
    \item \textbf{Couche métier~:} vues Django, services
    (\texttt{apps/chatbot/services.py}, \texttt{apps/prediction/services.py}),
    modèles ORM.
    \item \textbf{Couche données~:} PostgreSQL pour les entités relationnelles,
    système de fichiers pour les créatifs, fichier \texttt{model.pkl} pour
    le modèle ML pré-entraîné.
\end{itemize}

\subsection{Diagramme de cas d'utilisation}

Trois acteurs (Marketing Manager, Analyste, Administrateur) interagissent
avec dix cas d'utilisation principaux. Le cas \emph{«~Chat with Marketing
Assistant~»} est le différenciateur central du projet. Voir
\texttt{docs/diagrams.md} \S1.

\subsection{Diagramme de classes}

Le modèle de données comporte 11 entités principales~: \texttt{User},
\texttt{Product}, \texttt{Campaign}, \texttt{Angle}, \texttt{Creative},
\texttt{Hook}, \texttt{Expense}, \texttt{Revenue}, \texttt{ChatSession},
\texttt{ChatMessage}, \texttt{PredictionLog}.

Les relations clés sont~:
\begin{itemize}
    \item \texttt{Product 1---* Campaign} --- un produit peut faire l'objet
    de plusieurs campagnes successives.
    \item \texttt{Campaign 1---* Creative / Expense / Revenue}
    --- agrégations financières et créatives à la maille campagne.
    \item \texttt{Angle *---* Creative} et \texttt{Angle *---* Hook} ---
    réutilisation transversale des angles psychologiques.
    \item \texttt{ChatSession 1---* ChatMessage} --- historique
    conversationnel.
\end{itemize}

\subsection{Diagramme de séquence --- chatbot DeepSeek}

C'est le diagramme le plus important du projet car il illustre le mécanisme
de \emph{contextualisation} qui distingue notre chatbot d'un simple
\texttt{ChatGPT} générique.

\begin{enumerate}
    \item L'utilisateur saisit sa question dans le widget Alpine.js.
    \item Le navigateur envoie une requête \texttt{POST /chat/} au backend.
    \item La vue Django persiste le message utilisateur, puis appelle
    \texttt{build\_context(user)}.
    \item Le \emph{Context Builder} interroge la base pour extraire les
    statistiques pertinentes (top 5~campagnes, ROI moyen, meilleur angle
    du mois).
    \item Ces faits sont injectés dans le \texttt{system\_prompt} envoyé
    à l'API DeepSeek.
    \item DeepSeek retourne une complétion en streaming.
    \item Le backend persiste la réponse et la transmet au navigateur via
    \texttt{StreamingHttpResponse} (Server-Sent Events).
\end{enumerate}

Voir \texttt{docs/diagrams.md} \S3 pour le diagramme complet.

\subsection{Diagramme de séquence --- prédiction ML}

Le flux de prédiction est volontairement simple~: la vue
\texttt{/predict/} reçoit le formulaire, délègue à
\texttt{prediction.services.predict\_potential()}, qui charge \emph{une
seule fois} le modèle (mise en cache au démarrage de l'app) et retourne la
classe et la confiance. Une trace est écrite en base
(\texttt{PredictionLog}) pour permettre un ré-entraînement futur.

Voir \texttt{docs/diagrams.md} \S4.

\subsection{Diagramme d'états --- cycle de vie d'un produit}

Le produit transite par trois états~: \emph{Planning}~$\rightarrow$~%
\emph{Active}~$\rightarrow$~\emph{Évalué}, avec des transitions de retour
possibles (re-test avec un nouvel angle). Voir \texttt{docs/diagrams.md}
\S6.
